%2multibyte Version: 5.50.0.2953 CodePage: 936

\documentclass{article}
%%%%%%%%%%%%%%%%%%%%%%%%%%%%%%%%%%%%%%%%%%%%%%%%%%%%%%%%%%%%%%%%%%%%%%%%%%%%%%%%%%%%%%%%%%%%%%%%%%%%%%%%%%%%%%%%%%%%%%%%%%%%%%%%%%%%%%%%%%%%%%%%%%%%%%%%%%%%%%%%%%%%%%%%%%%%%%%%%%%%%%%%%%%%%%%%%%%%%%%%%%%%%%%%%%%%%%%%%%%%%%%%%%%%%%%%%%%%%%%%%%%%%%%%%%%%
\usepackage{amsmath}
\usepackage{amsfonts}

\setcounter{MaxMatrixCols}{10}
%TCIDATA{OutputFilter=LATEX.DLL}
%TCIDATA{Version=5.50.0.2953}
%TCIDATA{Codepage=936}
%TCIDATA{<META NAME="SaveForMode" CONTENT="1">}
%TCIDATA{BibliographyScheme=Manual}
%TCIDATA{Created=Monday, November 08, 2021 17:14:57}
%TCIDATA{LastRevised=Tuesday, November 08, 2022 22:16:12}
%TCIDATA{<META NAME="GraphicsSave" CONTENT="32">}
%TCIDATA{<META NAME="DocumentShell" CONTENT="Standard LaTeX\Blank - Standard LaTeX Article">}
%TCIDATA{CSTFile=40 LaTeX article.cst}

\newtheorem{theorem}{Theorem}
\newtheorem{acknowledgement}[theorem]{Acknowledgement}
\newtheorem{algorithm}[theorem]{Algorithm}
\newtheorem{axiom}[theorem]{Axiom}
\newtheorem{case}[theorem]{Case}
\newtheorem{claim}[theorem]{Claim}
\newtheorem{conclusion}[theorem]{Conclusion}
\newtheorem{condition}[theorem]{Condition}
\newtheorem{conjecture}[theorem]{Conjecture}
\newtheorem{corollary}[theorem]{Corollary}
\newtheorem{criterion}[theorem]{Criterion}
\newtheorem{definition}[theorem]{Definition}
\newtheorem{example}[theorem]{Example}
\newtheorem{exercise}[theorem]{Exercise}
\newtheorem{lemma}[theorem]{Lemma}
\newtheorem{notation}[theorem]{Notation}
\newtheorem{problem}[theorem]{Problem}
\newtheorem{proposition}[theorem]{Proposition}
\newtheorem{remark}[theorem]{Remark}
\newtheorem{solution}[theorem]{Solution}
\newtheorem{summary}[theorem]{Summary}
\newenvironment{proof}[1][Proof]{\noindent\textbf{#1.} }{\ \rule{0.5em}{0.5em}}
\input{tcilatex}
\begin{document}


\begin{center}
{\Large Class Test}
\end{center}

EACH QUESTION\ EQUALLY\ MARKED\ AT\ 12.5 MARKS\bigskip

\textbf{Question 1}

(a) Define an unbiased estimator and explain why this is an important
property. Explain how the Mean Squared Error (MSE)\ of an estimator is
related to its bias.

Solution:

An estimator $\hat{\theta}_{n}$ is called unbiased if $\mathbb{E}\hat{\theta}%
_{n}\mathbb{=\theta }$. In the case of unbiased estimators, $MSE\left( \hat{%
\theta}_{n},\theta \right) =var\left( \hat{\theta}_{n}\right) ,$ since the
bias $\mathbb{E}\hat{\theta}_{n}-\theta =0$. Otherwise, $MSE\left( \hat{%
\theta}_{n},\theta \right) =var\left( \hat{\theta}_{n}\right) +\left( 
\mathbb{E}\hat{\theta}_{n}-\theta \right) \left( \mathbb{E}\hat{\theta}%
_{n}-\theta \right) ^{\prime }.$

Award: 5 marks for precise definition of unbiasedness, 5 marks for precise
relation of bias to MSE, and 2.5 points for explanation, intuition.

(b) State the Central Limit Theorem for a sequence of i.i.d. random
variables carefully stating any conditions required. Briefly explain the
theorem's importance for statistical inference.

\bigskip

Solution:Let $X_{1},...,X_{n}$ be i.i.d. with $\mathbb{E}\left( X_{1}\right)
=\mu $ and $V\left( X_{1}\right) =\sigma ^{2}<\infty .$ Then, 
\begin{eqnarray*}
Z_{n} &=&\frac{\overline{X}_{n}-\mu }{\sqrt{V(\overline{X}_{n})}}=\frac{%
\sqrt{n}}{\sigma }(\overline{X}_{n}-\mu )=\frac{1}{\sigma \sqrt{n}}%
\sum_{i=1}^{n}(X_{i}-\mu )\rightarrow _{d}Z\text{ } \\
\text{\ where }Z &\sim &\mathcal{N}(0,1)
\end{eqnarray*}

\begin{theorem}
In other words, 
\begin{equation*}
\lim_{n\rightarrow \infty }F_{Z_{n}}=\Phi (z)=\int_{-\infty }^{z}(2\pi
)^{-1/2}\exp (-x^{2}/2)dx.
\end{equation*}
\end{theorem}

The CLT\ says that probability of $\frac{\sqrt{n}\left( \bar{X}_{n}-\mu
\right) }{\sigma }$ converges to a $N\left( 0,1\right) $\ r.v. This is a
remarkable result since nothing was assumed about the distribution of the
i.i.d. sequence $\left\{ X_{1},..,X_{n}\right\} $, apart from the existence
of their mean and variance; so the $X^{\prime }s$ could be generated by any
distribution. The theorem can be used for confidence interval construction
and hypothesis testing.

Award: 5 marks for statement + 4 marks for correct conditions; 3.5 marks for
explanation

(c) Explain what the difference is between an unbiased and a consistent
estimator. Does one imply the other? Explain why these are both desirable
properties for estimators to have.

Solution: \ For an unbiased estimator, $\mathbb{E}\left[ \hat{\theta}_{n}%
\right] =\theta ,$ while for a consistent estimator $\hat{\theta}%
_{n}\rightarrow \theta $ as $n\rightarrow \infty .$ The former is a finite
sample property ensuring that an estimator is not making systematic errors,
and tells us nothing about the limiting behaviour of $\hat{\theta}_{n}.$ The
latter is an asymptotic result that guarantees that as the sample size
increases, the distribution of $\hat{\theta}_{n}$ collapses to a point, the
true unknown $\theta .$ Roughly speaking consistency implies asymptotic
unbiasedness (whenever the expectation exists) and also implies that the
variance (when it exists) of the estimator is converging to zero (so the
precision increases as more information arrives).

Award: 8 marks for correct definition + 4.5 marks for explanation\bigskip

(d) Suppose $\sqrt{n}\left( \hat{\theta}_{n}-\theta \right) \rightarrow _{d}%
\mathcal{N}\left( 0,1\right) $ as $n\rightarrow \infty .$ Derive the
asymptotic distribution of $e^{2\widehat{\theta }_{n}}.$

\bigskip

Solution: Applying the delta method, we have that 
\begin{equation*}
\sqrt{n}\left( e^{2\hat{\theta}_{n}}-e^{2\theta }\right) \rightarrow _{d}%
\mathcal{N}\left( 0,4e^{4\theta }\right) .
\end{equation*}

Award: 5 point for correct approach, 5 for correct expression and 2.5 for
correct explanation. \bigskip

\textbf{Question 2}

Consider the linear regression model $y=X\beta +\varepsilon ,$ where $y$ is
an $n\times 1$ vector of dependent variables, $X$ is an $n\times k$ matrix
of regressors, $\beta $ is a $k\times 1$ vector of unknown parameters, and $%
\varepsilon $ is an $n\times 1$ vector of unobserved random variables. The
standard assumptions provided in the lecture are assumed to hold.

\begin{enumerate}
\item The linear regression model in is correctly specified and linear in $%
\beta $ and $\varepsilon .$

\item The columns of $X$ are linearly independent, so that $rank(X^{\prime
}X)=rank(X)=k<n.$

\item $\varepsilon _{i}$ are i.i.d. variables and $X$ are exogenous
identically distributed regressors, such that $\mathbb{E}[\varepsilon
_{i}|X]=0$ and $\mathbb{E}[\varepsilon _{i}^{2}|X]=\sigma ^{2}.$

\item $\limfunc{plim}_{n\rightarrow \infty }(\frac{1}{n}X^{\prime }X)=Q,$
where $\left\Vert Q\right\Vert <\infty $ and $Q$ is positive definite matrix.
\end{enumerate}

(a) Suppose there are two explanatory variables, $k=2,$ so that $%
y=X_{1}\beta _{1}+X_{2}\beta _{2}+\varepsilon .$ State without proof the
Frisch-Waugh-Lovell theorem for the OLS estimator $\hat{\beta}_{1}.$ Explain
how the theorem is related to omitted variable bias.

Solution:

\begin{theorem}[Frisch-Waugh-Lovell]
Suppose that we have a linear regression model and we take an arbitrary
partition $X=[X_{1},X_{2}]$ with $X_{1}$ an $n\times k_{1}$ and $X_{2}$ an $%
n\times \left( k-k_{1}\right) ,$ so that $y=X_{1}\beta _{1}+X_{2}\beta
_{2}+\varepsilon .$ Then, $\widehat{\beta }_{1}^{OLS}=(X_{1}^{\prime
}M_{X_{2}}X_{1})^{-1}X_{1}^{\prime }M_{X_{2}}y$ where $M_{X_{2}}$ is the
projection matrix defined as $I_{n}-X_{2}(X_{2}^{\prime
}X_{2})^{-1}X_{2}^{\prime }$.
\end{theorem}

Award: 6 for correct statement, 2 for omitted variable bias definiton and
4.5 for explanation.

(b) Given that the model is $y=X_{1}\beta _{1}+X_{2}\beta _{2}+\varepsilon ,$
show that the estimator for $\beta _{1}$ from Frisch-Waugh-Lovell theorem is
unbiased and derive its variance.

\begin{eqnarray*}
\hat{\beta}_{1} &=&\left( X_{1}^{\prime }M_{X_{2}}X_{1}\right)
^{-1}X_{1}^{\prime }M_{X_{2}}y=\left( X_{1}^{\prime }M_{X_{2}}X_{1}\right)
^{-1}X_{1}^{\prime }M_{X_{2}}\left( X_{1}\beta _{1}+X_{2}\beta
_{2}+\varepsilon \right) \\
&=&\beta _{1}+\left( X_{1}^{\prime }M_{X_{2}}X_{1}\right) ^{-1}X_{1}^{\prime
}M_{X_{2}}X_{2}\beta _{2}+\left( X_{1}^{\prime }M_{X_{2}}X_{1}\right)
^{-1}X_{1}^{\prime }M_{X_{2}}\varepsilon \\
&=&\beta _{1}+\left( X_{1}^{\prime }M_{X_{2}}X_{1}\right) ^{-1}X_{1}^{\prime
}M_{X_{2}}\varepsilon
\end{eqnarray*}%
since $M_{X_{2}}$ is orthogonal to $X_{2}.$ Taking expectations, we have 
\begin{equation*}
\mathbb{E}\left[ \hat{\beta}_{1}\right] =\beta _{1}+\mathbb{E}\left[ \left(
X_{1}^{\prime }M_{X_{2}}X_{1}\right) ^{-1}X_{1}^{\prime }M_{X_{2}}\mathbb{E}%
\left[ \varepsilon |X\right] \right] =\beta _{1}.
\end{equation*}
\begin{eqnarray*}
V\left[ \hat{\beta}_{1}\right] &=&\mathbb{E}\left[ \hat{\beta}_{1}-\mathbb{E}%
\hat{\beta}_{1}\right] \left[ \hat{\beta}_{1}-\mathbb{E}\hat{\beta}_{1}%
\right] ^{\prime } \\
&=&\mathbb{E}\left[ \left( X_{1}^{\prime }M_{X_{2}}X_{1}\right)
^{-1}X_{1}^{\prime }M_{X_{2}}\varepsilon \varepsilon ^{\prime
}M_{X_{2}}^{\prime }X_{1}^{\prime }\left( X_{1}^{\prime
}M_{X_{2}}X_{1}\right) ^{-1}\right] \\
&=&\mathbb{E}\left[ \left( X_{1}^{\prime }M_{X_{2}}X_{1}\right)
^{-1}X_{1}^{\prime }M_{X_{2}}\left[ \varepsilon \varepsilon ^{\prime
}|X_{1},X_{2}\right] M_{X_{2}}^{\prime }X_{1}^{\prime }\left( X_{1}^{\prime
}M_{X_{2}}X_{1}\right) ^{-1}\right] \\
&=&\sigma ^{2}\mathbb{E}\left[ \left( X_{1}^{\prime }M_{X_{2}}X_{1}\right)
^{-1}\right] .
\end{eqnarray*}

Award: 4 points for precise unbiasedness derivations, 4 for variance and 4.5
for explaning steps, using assumptions etc.

(c) Suppose you know that $\beta _{1}=\beta _{2}=0,$ so that the model
reduces to $y=\varepsilon .$ Is $\hat{\sigma}_{n}^{2}=\frac{1}{n}\varepsilon
^{\prime }\varepsilon $ unbiased? Derive an expression for the variance of $%
\hat{\sigma}_{n}^{2}.$

Solution: $\hat{\sigma}^{2}=\frac{1}{n}\varepsilon ^{\prime }\varepsilon $
as the $\varepsilon $ is observed and $\mathbb{E}\left[ \hat{\sigma}^{2}%
\right] =\mathbb{E}\left[ \frac{1}{n}\tsum\nolimits_{i=1}^{n}\varepsilon
_{i}^{2}\right] =\sigma ^{2}.$ Alternatively, since $\varepsilon $ is
observed%
\begin{gather*}
\mathbb{E}[\widehat{\sigma }_{n}^{2}]=\mathbb{E}[\frac{1}{n}\widehat{%
\varepsilon }^{\prime }\widehat{\varepsilon }]=\mathbb{E}[\frac{1}{n}%
\varepsilon ^{\prime }\varepsilon ]=\frac{1}{n}\mathbb{E}[tr(\varepsilon
\varepsilon ^{\prime })] \\
=\frac{1}{n}tr(\mathbb{E}[\varepsilon \varepsilon ^{\prime }])=\frac{\sigma
^{2}}{n}\mathbb{E[}tr(I_{n})]=\sigma ^{2}\text{.}
\end{gather*}%
Let $z_{i}=\varepsilon _{i}^{2}-\sigma ^{2}$ which has expectation 0, 
\begin{gather*}
V\left( \widehat{\sigma }_{n}^{2}\right) =\mathbb{E}\left( \frac{1}{n}%
\tsum\nolimits_{i=1}^{n}\varepsilon _{i}^{2}-\sigma ^{2}\right) ^{2}= \\
=\mathbb{E}\left( \frac{1}{n}\tsum\nolimits_{i=1}^{n}\left( \varepsilon
_{i}^{2}-\sigma ^{2}\right) \right) ^{2} \\
=\frac{1}{n^{2}}\mathbb{E}\left( \tsum\nolimits_{i=1}^{n}z_{i}\right) ^{2}%
\overset{\text{0 mean, ind}}{=}\frac{1}{n^{2}}\tsum\nolimits_{i=1}^{n}%
\mathbb{E}z_{i}^{2} \\
\overset{\text{i.d.}}{=}\frac{1}{n}\left( \mathbb{E}z_{1}^{2}\right) =\frac{1%
}{n}\mathbb{E}\left( \varepsilon _{1}^{2}-\sigma ^{2}\right) ^{2}=\frac{1}{n}%
\left( \mathbb{E}\varepsilon _{1}^{4}-\sigma ^{4}\right)
\end{gather*}
So we additionally require $\mathbb{E}\left( \varepsilon _{1}^{4}\right)
<\infty .$ $\hat{\sigma}_{n}^{2}$ is unbiased, this is because no additional
parameters need to be estimated.

Award: 5 point for unbiasedness proof with precision, 5 points for the
variance derivation; 2.5 for explaining all steps, intuition, requiring 4th
moment.

\textbf{Question 3 }

(a) Consider the linear regression model from Question 2. Recall that the
ridge estimator is defined as $\hat{\beta}^{Ridge}=\left( X^{\prime
}X+\lambda _{n}I_{k}\right) ^{-1}X^{\prime }y.$ Suppose that $\lambda
_{n}=cn^{\delta /2}$ for some $\delta \geq 0.$ For which range of values for 
$c$ and $\delta $ is $\hat{\beta}^{Ridge}$ a consistent estimator for $\beta
?$

Solution: 
\begin{eqnarray*}
p\lim_{n\rightarrow \infty }\hat{\beta}^{Ridge} &=&\left(
p\lim_{n\rightarrow \infty }\frac{1}{n}\tsum\nolimits_{i=1}^{n}\mathbf{x}_{i}%
\mathbf{x}_{i}^{\prime }+\lim_{n\rightarrow \infty }cn^{\delta
/2-1}I_{k}\right) ^{-1}p\lim_{n\rightarrow \infty }\frac{1}{n}%
\tsum\nolimits_{i=1}^{n}\mathbf{x}_{i}\mathbf{x}_{i}^{\prime }\beta \\
&&+\left( p\lim_{n\rightarrow \infty }\frac{1}{n}\tsum\nolimits_{i=1}^{n}%
\mathbf{x}_{i}\mathbf{x}_{i}^{\prime }+\lim_{n\rightarrow \infty }cn^{\delta
/2-1}I_{k}\right) ^{-1}p\lim_{n\rightarrow \infty }\frac{1}{n}%
\tsum\nolimits_{i=1}^{n}\mathbf{x}_{i}\mathbf{\varepsilon }_{i}
\end{eqnarray*}

so for any values of $c\in \left( -\infty ,\infty \right) $ and any $\delta
\in \lbrack 0,2),$ $\lim_{n\rightarrow \infty }cn^{\delta /2-1}I_{k}=0$ so $%
p\lim_{n\rightarrow \infty }\hat{\beta}^{Ridge}=\beta .$

Award: 7.5 point for correct derivation of plim, 5 points for correct values
of $c$ and $\delta .$

\end{document}
